\documentclass[10pt]{article}
\usepackage[a4paper,landscape,margin=1.6cm]{geometry}
\usepackage[table]{xcolor}
\usepackage{adjustbox}
\usepackage{booktabs,graphicx,amsmath,float,caption,enumitem,hyperref}
\captionsetup{font=small,labelfont=bf}
\setlength{\parindent}{0pt}\setlength{\parskip}{4pt}
\definecolor{good}{RGB}{190,232,190}
\definecolor{fair}{RGB}{252,226,160}
\definecolor{poor}{RGB}{244,182,176}
\definecolor{boxbg}{RGB}{238,242,248}
\newcommand{\fig}[3][0.9]{\begin{figure}[H]\centering\includegraphics[width=#1\linewidth]{figures/#2}\caption{#3}\end{figure}}
\newcommand{\what}[1]{\par\colorbox{boxbg}{\parbox{\dimexpr\linewidth-2\fboxsep}{\small #1}}\par}

\title{Amplifier and setup scorecard:\\PA03 (130\,K) and the LSM production modules DM07/DM09/DM10 (180\,K)}
\author{amp\_phys analysis}
\date{16 November 2026}

\begin{document}
\maketitle

\section*{What is compared, and how to read the colours}
Every number here comes from the \textbf{raw readout traces} (plus, where marked, the images of the same channel), measured
the same way for every amplifier:
\begin{itemize}[nosep]
\item \textbf{PA03}, board 101 (ch0--ch3, four CCDs) and board 102 (every channel has an output transistor, but no charge is clocked to ch0--ch2, and none reaches ch3 in
  the traces: a \emph{same-board reference without charge}), November v3 data, 130\,K. Trace numbers are medians over 8 traces; the
  ``p90'' columns use all 93 packets whose images are valid. Operating points where the readout breaks (merged peaks, most columns hot; flagged by the image fit) are left out everywhere: their traces do not describe a working amplifier, and they are a separate valid/invalid outcome for the optimiser.
\item \textbf{LSM} production modules DM07/DM09/DM10 on ACM boards 106, 109 and 110, 180\,K, one trace and one static
  noise record per channel, and the 500/1000-skip images processed by the LSM pipeline (gain, noise, dark charge).
\end{itemize}
Each graded cell is \colorbox{good}{good}, \colorbox{fair}{fair} or \colorbox{poor}{poor} against the thresholds printed
under its table. The thresholds are physics or practice limits (e.g.\ good thermal oxide, the frontend noise floor), not
rankings, so a whole column can be green or red. Uncoloured columns are context, not graded.
Table~\ref{tab:card} condenses everything into one card; the following sections explain each quantity, how it is
measured and why it matters.

\begin{table}[H]\centering\small
\caption{Report card: the typical grade of each category per channel (detail in Tables~\ref{tab:sn}--\ref{tab:setup}).}\label{tab:card}
\begin{adjustbox}{max width=\linewidth}
\begin{tabular}{llcccccc}\toprule
setup & channel & Sense node & Read noise & Transistor (1/f, traps) & Timing & Setup & Charge (vs same board)\\\midrule
PA03, 130 K & 101 ch0 & \cellcolor{good}good \scriptsize{1/0/0} & \cellcolor{good}good \scriptsize{4/0/0} & \cellcolor{good}good \scriptsize{4/0/0} & \cellcolor{good}good \scriptsize{3/0/0} & \cellcolor{poor}poor \scriptsize{0/1/2} & \cellcolor{good}good \scriptsize{1/0/0}\\
 & 101 ch1 & \cellcolor{good}good \scriptsize{1/0/0} & \cellcolor{good}good \scriptsize{3/1/0} & \cellcolor{good}good \scriptsize{3/1/0} & \cellcolor{good}good \scriptsize{3/0/0} & \cellcolor{poor}poor \scriptsize{0/1/2} & \cellcolor{good}good \scriptsize{1/0/0}\\
 & 101 ch2 & \cellcolor{good}good \scriptsize{1/0/0} & \cellcolor{fair}fair \scriptsize{2/2/0} & \cellcolor{fair}fair \scriptsize{2/2/0} & \cellcolor{good}good \scriptsize{3/0/0} & \cellcolor{poor}poor \scriptsize{0/1/2} & \cellcolor{fair}fair \scriptsize{0/1/0}\\
 & 101 ch3 & \cellcolor{good}good \scriptsize{1/0/0} & \cellcolor{fair}fair \scriptsize{2/2/0} & \cellcolor{good}good \scriptsize{3/1/0} & \cellcolor{good}good \scriptsize{3/0/0} & \cellcolor{poor}poor \scriptsize{0/1/2} & \cellcolor{good}good \scriptsize{1/0/0}\\
 & 102 ch3 & \cellcolor{good}good \scriptsize{1/0/0} & \cellcolor{good}good \scriptsize{3/1/0} & \cellcolor{good}good \scriptsize{3/1/0} & \cellcolor{fair}fair \scriptsize{1/2/0} & \cellcolor{good}good \scriptsize{3/0/0} & \cellcolor{good}good \scriptsize{1/0/0}\\
 & 102 ch0 \textit{(no charge)} & -- & \cellcolor{good}good \scriptsize{1/0/0} & \cellcolor{good}good \scriptsize{2/0/0} & \cellcolor{fair}fair \scriptsize{0/1/0} & \cellcolor{good}good \scriptsize{3/0/0} & --\\
 & 102 ch1 \textit{(no charge)} & -- & -- & \cellcolor{good}good \scriptsize{1/0/0} & -- & -- & --\\
 & 102 ch2 \textit{(no charge)} & -- & -- & \cellcolor{good}good \scriptsize{1/0/0} & -- & -- & --\\
\midrule
LSM, 180 K & ACM\,106 ch0 & \cellcolor{good}good \scriptsize{1/0/0} & \cellcolor{fair}fair \scriptsize{2/2/0} & \cellcolor{good}good \scriptsize{2/1/0} & \cellcolor{good}good \scriptsize{3/0/0} & \cellcolor{fair}fair \scriptsize{2/2/0} & \cellcolor{good}good \scriptsize{1/0/0}\\
 & ACM\,106 ch1 & \cellcolor{good}good \scriptsize{1/0/0} & \cellcolor{fair}fair \scriptsize{2/1/1} & \cellcolor{good}good \scriptsize{2/1/0} & \cellcolor{good}good \scriptsize{2/1/0} & \cellcolor{good}good \scriptsize{3/1/0} & \cellcolor{good}good \scriptsize{1/0/0}\\
 & ACM\,106 ch2 & \cellcolor{good}good \scriptsize{1/0/0} & \cellcolor{fair}fair \scriptsize{2/2/0} & \cellcolor{good}good \scriptsize{2/1/0} & \cellcolor{good}good \scriptsize{2/1/0} & \cellcolor{fair}fair \scriptsize{2/2/0} & \cellcolor{good}good \scriptsize{1/0/0}\\
 & ACM\,106 ch3 & \cellcolor{good}good \scriptsize{1/0/0} & \cellcolor{fair}fair \scriptsize{2/2/0} & \cellcolor{good}good \scriptsize{2/1/0} & \cellcolor{good}good \scriptsize{2/1/0} & \cellcolor{good}good \scriptsize{3/1/0} & \cellcolor{good}good \scriptsize{1/0/0}\\
 & ACM\,109 ch0 & \cellcolor{good}good \scriptsize{1/0/0} & \cellcolor{fair}fair \scriptsize{2/2/0} & \cellcolor{good}good \scriptsize{2/1/0} & \cellcolor{good}good \scriptsize{2/1/0} & \cellcolor{fair}fair \scriptsize{2/2/0} & \cellcolor{good}good \scriptsize{1/0/0}\\
 & ACM\,109 ch1 & \cellcolor{good}good \scriptsize{1/0/0} & \cellcolor{fair}fair \scriptsize{0/4/0} & \cellcolor{poor}poor \scriptsize{1/0/2} & \cellcolor{fair}fair \scriptsize{1/1/1} & \cellcolor{poor}poor \scriptsize{1/1/2} & \cellcolor{good}good \scriptsize{1/0/0}\\
 & ACM\,109 ch2 & \cellcolor{good}good \scriptsize{1/0/0} & \cellcolor{fair}fair \scriptsize{2/2/0} & \cellcolor{good}good \scriptsize{2/1/0} & \cellcolor{fair}fair \scriptsize{1/1/1} & \cellcolor{good}good \scriptsize{3/0/1} & \cellcolor{good}good \scriptsize{1/0/0}\\
 & ACM\,109 ch3 & \cellcolor{good}good \scriptsize{1/0/0} & \cellcolor{fair}fair \scriptsize{2/2/0} & \cellcolor{good}good \scriptsize{1/0/0} & \cellcolor{fair}fair \scriptsize{1/1/1} & \cellcolor{good}good \scriptsize{3/1/0} & \cellcolor{good}good \scriptsize{1/0/0}\\
 & ACM\,110 ch0 & \cellcolor{good}good \scriptsize{1/0/0} & \cellcolor{good}good \scriptsize{4/0/0} & \cellcolor{good}good \scriptsize{2/1/0} & \cellcolor{good}good \scriptsize{2/1/0} & \cellcolor{good}good \scriptsize{4/0/0} & \cellcolor{good}good \scriptsize{1/0/0}\\
 & ACM\,110 ch1 & \cellcolor{good}good \scriptsize{1/0/0} & \cellcolor{fair}fair \scriptsize{2/2/0} & \cellcolor{good}good \scriptsize{2/0/1} & \cellcolor{good}good \scriptsize{2/0/1} & \cellcolor{good}good \scriptsize{4/0/0} & \cellcolor{good}good \scriptsize{1/0/0}\\
 & ACM\,110 ch2 & \cellcolor{good}good \scriptsize{1/0/0} & \cellcolor{good}good \scriptsize{4/0/0} & \cellcolor{good}good \scriptsize{2/1/0} & \cellcolor{good}good \scriptsize{2/0/1} & \cellcolor{good}good \scriptsize{4/0/0} & \cellcolor{good}good \scriptsize{1/0/0}\\
 & ACM\,110 ch3 & \cellcolor{good}good \scriptsize{1/0/0} & \cellcolor{fair}fair \scriptsize{1/3/0} & \cellcolor{fair}fair \scriptsize{1/2/0} & \cellcolor{good}good \scriptsize{2/0/1} & \cellcolor{poor}poor \scriptsize{1/1/2} & \cellcolor{good}good \scriptsize{1/0/0}\\
\bottomrule\end{tabular}\end{adjustbox}
\par\smallskip\footnotesize Colour = the typical (median) grade of the category; the small numbers count its graded quantities as good/fair/poor, so a single problem stays visible without colouring the whole category; ``--'': not measurable for that channel.
\end{table}


\paragraph{The card in one paragraph.} The PA03 amplifiers have low 1/f corners; 101 ch2's noise is mildly heavy-tailed (its noise level is modulated).
Board 101 carries the most slow pickup measured ($\sim$35\,\textmu V/$\sqrt{\text{Hz}}$), but its 500-skip images still reach
the white-noise expectation, so the CDS and skip averaging reject most of it. The LSM setup is much quieter in its
traces and its transistors have similar 1/f noise, but its
\textbf{images are noisier} at 180\,K and its \textbf{timing after the summing well is not settled} on several channels.
Most channels are good or fair in most categories; all 12 LSM amplifiers work (every one resolves single electrons, 0.16--0.23\,e at 1000 skips); the channels with a
noise feature to follow up are ACM109 ch1 (intermittent $\sim$1\,ms noise bursts, one trace), ACM110 ch3 (interference), and on PA03 the setup of board 101 and 101 ch2's extra charge and telegraph
episodes (Section~\ref{sec:verdicts}).

% ===================================================================================================================
\clearpage
\section{Sense node: how many microvolts one electron makes}
\what{\textbf{Conversion gain} $CG=\text{gain}/G$: the image gain (ADU per electron, from the single-electron peaks)
divided by the readout-chain gain $G$ (0.875\,ADU/\textmu V by design; the earlier 0.853 from board 102 assumed an open input, which does not hold,
since those channels have transistors). \textbf{Sense-node capacitance} $C_\text{sn}=A\,q/CG$ with follower gain $A=0.85$ (layout).
\textbf{Why it matters:} a smaller node gives more volts per electron, so every voltage noise source costs fewer
electrons. Everything else being equal, the read noise in electrons scales with $C_\text{sn}$.}
\begin{table}[H]\centering\small
\caption{Sense node: conversion gain and capacitance ($A=0.85$).}\label{tab:sn}
\begin{adjustbox}{max width=\linewidth}
\begin{tabular}{llrrr}\toprule
setup & channel & gain (ADU/e) & conv.\ gain (\textmu V/e) & $C_\text{sn}$ (fF)\\\midrule
PA03, 130 K & 101 ch0 & 1.41 & 1.62 & \cellcolor{good}84\\
 & 101 ch1 & 1.39 & 1.59 & \cellcolor{good}85\\
 & 101 ch2 & 1.33 & 1.52 & \cellcolor{good}89\\
 & 101 ch3 & 1.34 & 1.53 & \cellcolor{good}89\\
 & 102 ch3 & 1.29 & 1.51 & \cellcolor{good}90\\
 & 102 ch0 \textit{(no charge)} & -- & -- & --\\
 & 102 ch1 \textit{(no charge)} & -- & -- & --\\
 & 102 ch2 \textit{(no charge)} & -- & -- & --\\
\midrule
LSM, 180 K & ACM\,106 ch0 & 1.26 & 1.44 & \cellcolor{good}95\\
 & ACM\,106 ch1 & 1.27 & 1.45 & \cellcolor{good}94\\
 & ACM\,106 ch2 & 1.28 & 1.46 & \cellcolor{good}93\\
 & ACM\,106 ch3 & 1.31 & 1.49 & \cellcolor{good}91\\
 & ACM\,109 ch0 & 1.38 & 1.58 & \cellcolor{good}86\\
 & ACM\,109 ch1 & 1.37 & 1.57 & \cellcolor{good}87\\
 & ACM\,109 ch2 & 1.34 & 1.53 & \cellcolor{good}89\\
 & ACM\,109 ch3 & 1.31 & 1.49 & \cellcolor{good}91\\
 & ACM\,110 ch0 & 1.31 & 1.50 & \cellcolor{good}91\\
 & ACM\,110 ch1 & 1.33 & 1.52 & \cellcolor{good}89\\
 & ACM\,110 ch2 & 1.37 & 1.56 & \cellcolor{good}87\\
 & ACM\,110 ch3 & 1.28 & 1.46 & \cellcolor{good}93\\
\bottomrule\end{tabular}\end{adjustbox}
\par\smallskip\footnotesize Grades --- $C_\text{sn}$: \colorbox{good}{$x\le95$} \colorbox{fair}{$x\le105$}; otherwise \colorbox{poor}{poor}.
\end{table}

\begin{itemize}[nosep]
\item All 17 amplifiers are the same design: 84--95\,fF. The LSM nodes are $\sim$5\,\% larger on average (lower gain),
  with a 10\,\% spread between channels; PA03 101 ch0/ch1 have the smallest nodes.
\item The absolute $C_\text{sn}$ scales with the assumed $A$ and $G$ (both $\pm$5\,\%); the comparison between channels does not.
\end{itemize}

% ===================================================================================================================
\section{Read noise}
\what{\textbf{CDS noise}: the white noise of one skip (pedestal minus signal, 10\,\textmu s windows), converted to the
video input: $\sigma_1/G\cdot\sqrt{\tau}$ in nV/$\sqrt{\text{Hz}}$. An open input (load resistor + preamp only) would give 17.1; the board-102 channels without charge measure 17.0
(transistor present: a reference, not the bare frontend). \textbf{Per skip (e)}: the same in electrons ($\sigma_1/\text{gain}$). \textbf{Image}: the read noise the
images reach after 500 skips (zero-electron peak width). \textbf{Image/expectation}: that, divided by white-noise averaging
$\sigma_1/\sqrt{500}$: 1 = the skips average perfectly; above 1, slow noise or charge arriving during the skipping adds.}
\begin{table}[H]\centering\small
\caption{Read noise: per skip at the input and in electrons, and what the images reach after 500 skips.}\label{tab:noise}
\begin{adjustbox}{max width=\linewidth}
\begin{tabular}{llrrrr}\toprule
setup & channel & CDS (nV/$\sqrt{\text{Hz}}$) & per skip (e) & image, 500 skips (e) & image / expectation\\\midrule
PA03, 130 K & 101 ch0 & \cellcolor{fair}18.7 & \cellcolor{good}3.67 & \cellcolor{good}0.144 & \cellcolor{good}0.88\\
 & 101 ch1 & \cellcolor{fair}19.6 & \cellcolor{fair}3.90 & \cellcolor{good}0.147 & \cellcolor{good}0.84\\
 & 101 ch2 & \cellcolor{poor}20.4 & \cellcolor{poor}4.24 & \cellcolor{fair}0.164 & \cellcolor{good}0.86\\
 & 101 ch3 & \cellcolor{poor}21.1 & \cellcolor{poor}4.36 & \cellcolor{fair}0.173 & \cellcolor{good}0.89\\
 & 102 ch3 & \cellcolor{fair}18.7 & \cellcolor{fair}4.02 & \cellcolor{fair}0.200 & \cellcolor{fair}1.11\\
 & 102 ch0 \textit{(dead)} & \cellcolor{good}17.0 & -- & -- & --\\
\midrule
LSM, 180 K & ACM\,106 ch0 & \cellcolor{good}16.4 & \cellcolor{good}3.62 & \cellcolor{poor}0.225 & \cellcolor{poor}1.39\\
 & ACM\,106 ch1 & \cellcolor{good}16.9 & \cellcolor{good}3.68 & \cellcolor{poor}0.328 & \cellcolor{poor}1.99\\
 & ACM\,106 ch2 & \cellcolor{good}16.9 & \cellcolor{good}3.66 & \cellcolor{poor}0.208 & \cellcolor{fair}1.27\\
 & ACM\,106 ch3 & \cellcolor{good}17.1 & \cellcolor{good}3.61 & \cellcolor{poor}0.274 & \cellcolor{poor}1.70\\
 & ACM\,109 ch0 & \cellcolor{good}16.6 & \cellcolor{good}3.31 & \cellcolor{poor}0.262 & \cellcolor{poor}1.77\\
 & ACM\,109 ch1 & \cellcolor{poor}20.4 & \cellcolor{fair}4.12 & \cellcolor{poor}0.273 & \cellcolor{poor}1.48\\
 & ACM\,109 ch2 & \cellcolor{good}17.3 & \cellcolor{good}3.59 & \cellcolor{poor}0.226 & \cellcolor{poor}1.41\\
 & ACM\,109 ch3 & \cellcolor{fair}17.5 & \cellcolor{good}3.71 & \cellcolor{poor}0.215 & \cellcolor{fair}1.29\\
 & ACM\,110 ch0 & \cellcolor{good}17.3 & \cellcolor{good}3.66 & \cellcolor{fair}0.199 & \cellcolor{fair}1.22\\
 & ACM\,110 ch1 & \cellcolor{poor}20.9 & \cellcolor{poor}4.33 & \cellcolor{fair}0.193 & \cellcolor{good}1.00\\
 & ACM\,110 ch2 & \cellcolor{fair}17.8 & \cellcolor{good}3.60 & \cellcolor{fair}0.192 & \cellcolor{fair}1.19\\
 & ACM\,110 ch3 & \cellcolor{fair}19.7 & \cellcolor{poor}4.27 & \cellcolor{poor}0.209 & \cellcolor{fair}1.09\\
\bottomrule\end{tabular}\end{adjustbox}
\par\smallskip\footnotesize Grades --- CDS noise: \colorbox{good}{$x\le17.5$} \colorbox{fair}{$x\le20$}; per skip: \colorbox{good}{$x\le3.8$} \colorbox{fair}{$x\le4.2$}; image: \colorbox{good}{$x\le0.16$} \colorbox{fair}{$x\le0.2$}; image/expectation: \colorbox{good}{$x\le1$} \colorbox{fair}{$x\le1.3$}; otherwise \colorbox{poor}{poor}.
\end{table}

\begin{itemize}[nosep]
\item Per skip, both setups are at 3.3--4.4\,e. The LSM traces are slightly quieter at the input (16.4--17.8\,nV/$\sqrt{\text{Hz}}$
  on the good channels, against 18.7--21.1), but this relies on our chain gain for their boards.
\item \textbf{The images invert the ranking.} PA03 reaches 0.144--0.173\,e at 500 skips (slightly \emph{better} than white
  averaging); the LSM images reach 0.19--0.33\,e, 1.0--2.0$\times$ worse than their own traces predict. At 180\,K the blank
  images carry $\sim$1\,e/pixel: charge arriving while a pixel is skipped broadens the zero-electron peak.
\end{itemize}

% ===================================================================================================================
\clearpage
\section{The output transistor: 1/f noise, oxide traps, telegraph noise}
\what{The CDS removes slow noise only partly; what is left of the transistor's 1/f noise sits at 10--40\,kHz and does not
average away (Figure~\ref{fig:cds}). White noise falls with a longer CDS window, 1/f does not. Forming windows of 2--9.3\,\textmu s
\emph{in software} from the samples of each 10\,\textmu s window, and fitting $\sigma_1^2(\tau)=S_wI_w(\tau)+K\,I_f(\tau)$
with the exact filter, gives the \textbf{white part}, the \textbf{1/f corner} $K/S_w$, the \textbf{1/f share} of the
10\,\textmu s noise, and through the McWhorter model the \textbf{oxide trap density}
$N_t=K_{V_g}WL\,C_\text{ox}^2/(q^2kT\lambda)$ (good thermal oxide: $10^{16}$--$10^{17}$\,eV$^{-1}$cm$^{-3}$). These are \emph{upper limits}: the fit has not been
run on an open input, and band-limited white noise can be counted as 1/f. Consistent with the hardware window scan (100 strips) on PA03: within 15\,\% on ch1/ch3, $\times$1.6 on ch0/ch2.
\textbf{Kurtosis} of the skip-to-skip differences: 0 for Gaussian noise, positive for telegraph (single-trap) noise.}
\fig[0.85]{cds_transfer.png}{Why the CDS does not remove 1/f noise (left) and why skip averaging cannot separate it from
white noise (right).\label{fig:cds}}
\begin{table}[H]\centering\small
\caption{Output transistor: white and 1/f noise, oxide trap density, telegraph noise.}\label{tab:tr}
\begin{adjustbox}{max width=\linewidth}
\begin{tabular}{llrrrrrrr}\toprule
setup & channel & white (nV/$\sqrt{\text{Hz}}$) & 1/f corner (kHz) & 1/f share & $N_t$ ($10^{16}$/eV/cm$^3$) & kurtosis & kurtosis p90 & fit resid.\\\midrule
PA03, 130 K & 101 ch0 & \cellcolor{fair}14.2 & \cellcolor{good}16 & 0.37 & \cellcolor{good}1.10 & \cellcolor{good}$-$0.02 & \cellcolor{good}0.07 & \cellcolor{good}0.026\\
 & 101 ch1 & \cellcolor{fair}12.6 & \cellcolor{good}34 & 0.56 & \cellcolor{good}1.84 & \cellcolor{good}0.01 & \cellcolor{good}0.10 & \cellcolor{good}0.019\\
 & 101 ch2 & \cellcolor{fair}14.1 & \cellcolor{good}27 & 0.50 & \cellcolor{good}1.78 & \cellcolor{good}$-$0.01 & \cellcolor{fair}0.17 & \cellcolor{good}0.020\\
 & 101 ch3 & \cellcolor{fair}13.0 & \cellcolor{good}39 & 0.59 & \cellcolor{fair}2.27 & \cellcolor{good}$-$0.01 & \cellcolor{good}0.09 & \cellcolor{good}0.009\\
 & 102 ch3 & \cellcolor{good}7.8 & \cellcolor{poor}123 & 0.82 & \cellcolor{fair}2.53 & \cellcolor{good}$-$0.06 & \cellcolor{good}0.10 & \cellcolor{good}0.011\\
 & 102 ch0 \textit{(dead)} & -- & -- & -- & -- & \cellcolor{good}$-$0.05 & -- & --\\
\midrule
LSM, 180 K & ACM\,106 ch0 & \cellcolor{good}7.5 & \cellcolor{fair}99 & 0.79 & \cellcolor{good}1.36 & \cellcolor{good}0.20 & -- & \cellcolor{good}0.019\\
 & ACM\,106 ch1 & \cellcolor{good}8.4 & \cellcolor{fair}75 & 0.74 & \cellcolor{good}1.30 & \cellcolor{good}0.09 & -- & \cellcolor{good}0.019\\
 & ACM\,106 ch2 & \cellcolor{good}9.1 & \cellcolor{fair}61 & 0.70 & \cellcolor{good}1.24 & \cellcolor{good}$-$0.02 & -- & \cellcolor{good}0.013\\
 & ACM\,106 ch3 & \cellcolor{good}8.4 & \cellcolor{fair}77 & 0.74 & \cellcolor{good}1.35 & \cellcolor{good}$-$0.08 & -- & \cellcolor{good}0.008\\
 & ACM\,109 ch0 & \cellcolor{good}8.6 & \cellcolor{fair}66 & 0.71 & \cellcolor{good}1.21 & \cellcolor{good}$-$0.16 & -- & \cellcolor{good}0.015\\
 & ACM\,109 ch1 & \cellcolor{good}5.4 & \cellcolor{poor}334 & 0.93 & \cellcolor{fair}2.40 & \cellcolor{poor}2.66 & -- & \cellcolor{good}0.012\\
 & ACM\,109 ch2 & \cellcolor{good}8.3 & \cellcolor{fair}83 & 0.76 & \cellcolor{good}1.43 & \cellcolor{good}$-$0.29 & -- & \cellcolor{good}0.019\\
 & ACM\,109 ch3 & -- & -- & -- & -- & \cellcolor{good}$-$0.23 & -- & \cellcolor{poor}0.210\\
 & ACM\,110 ch0 & \cellcolor{good}7.9 & \cellcolor{fair}93 & 0.78 & \cellcolor{good}1.43 & \cellcolor{good}$-$0.03 & -- & \cellcolor{good}0.014\\
 & ACM\,110 ch1 & \cellcolor{good}7.4 & \cellcolor{poor}174 & 0.87 & \cellcolor{fair}2.37 & \cellcolor{good}$-$0.02 & -- & \cellcolor{good}0.010\\
 & ACM\,110 ch2 & \cellcolor{good}8.6 & \cellcolor{fair}83 & 0.76 & \cellcolor{good}1.51 & \cellcolor{good}0.07 & -- & \cellcolor{good}0.010\\
 & ACM\,110 ch3 & \cellcolor{good}8.3 & \cellcolor{fair}117 & 0.81 & \cellcolor{good}1.97 & \cellcolor{fair}0.76 & -- & \cellcolor{good}0.009\\
\bottomrule\end{tabular}\end{adjustbox}
\par\smallskip\footnotesize Grades --- white: \colorbox{good}{$x\le10$} \colorbox{fair}{$x\le15$}; 1/f corner: \colorbox{good}{$x\le50$} \colorbox{fair}{$x\le120$}; $N_t$: \colorbox{good}{$x\le2$} \colorbox{fair}{$x\le5$}; kurtosis: \colorbox{good}{$|x|\le0.3$} \colorbox{fair}{$|x|\le1$}; kurtosis p90: \colorbox{good}{$|x|\le0.15$} \colorbox{fair}{$|x|\le0.5$}; otherwise \colorbox{poor}{poor}. Fit residual above 5\,\% $\Rightarrow$ the trap numbers are withheld.
\end{table}

\begin{itemize}[nosep]
\item \textbf{The health number: the transistor's 1/f noise left in a 500-skip pixel}, in electrons:
  $\sqrt{K\,I_f(500)}/\text{gain}$, with $K$ the fitted 1/f level and $I_f$ the exact CDS + 500-skip filter. It uses only
  the 1/f \emph{level} $K$, which is robust: on the LSM channels $K$ from the readout trace agrees with $K$ fitted
  directly to the static-record spectrum (no CDS model, clocks frozen) to $1.01\pm0.13$. The white/1/f \emph{split}
  and the corner frequency do not agree between the two (white 1.3$\times$ higher during readout), so they are shown but
  not graded. On PA03 board 101 the transistor's 1/f noise is 0.08--0.13\,e, i.e.\ 34--57\,\% of the image variance at
  500 skips: about half of what the images see comes from the transistor itself.
\item \textbf{Similar in both module sets:} $N_t\le1.1$--$2.3\times10^{16}$ (upper limits) on every good channel.
\item The LSM 1/f corners (61--99\,kHz) are higher than PA03's (16--39\,kHz): at 180\,K the 1/f power grows with $kT$
  and the white part is lower. Graded the same; read the LSM column with that in mind.
\item \textbf{What the heavy tails are.} Single-trap (telegraph) steps were searched for directly (two-level fits to the
  smoothed CDS): none is resolved. One trapped charge would move the output by only $\sim$0.76\,ADU against 5\,ADU per skip,
  and the fitted dwell times ($\sim$0.2\,ms) are too short to average it out. ACM109 ch1's kurtosis (2.7) comes from
  $\sim$1\,ms \emph{noise bursts} (pairs of pixels swinging by $\pm$20--30\,ADU), not a trap; its static spectrum is
  normal, so its high trace corner (334\,kHz) is the bursts. PA03 101 ch2's tails (kurtosis p90 0.17, maximum 0.70) come
  from a \emph{modulated noise level}: its per-pixel noise varies by 1.8\,\% beyond statistics (others 0.95--1.4\,\%), as a
  slowly switching fluctuator modulating the transistor noise would do. ACM110 ch1 has a higher 1/f level (174\,kHz corner).
\item ACM109 ch3: the fit fails (21\,\% residual, its window is still settling), so no trap numbers.
\end{itemize}
\fig[0.75]{oxide_traps_tau.png}{The hardware version of the method (PA03, 100 strips): noise against the CDS window,
with the white + 1/f fits.}

% ===================================================================================================================
\clearpage
\section{Timing and clocks}
\what{\textbf{Reset settling}: the time constant of the output after the reset/OG pulse (fit to the median skip cycle):
the output node's bandwidth. \textbf{Window slope}: the slope of the signal window in electrons per \textmu s: a window
that is still settling (it cancels on average, but jitter and drifts do not). \textbf{CDS offset}: signal minus pedestal
level of the empty-pixel waveform, the residual summing-well/output-gate transient. \textbf{RG, H, V}: the size of each
clock's waveform at the output per volt of swing, in electrons: \emph{relative} numbers that include board pickup (the
per-line mutual capacitances in fF need a voltage scan; PA03 only, Section~\ref{sec:scan}).}
\begin{table}[H]\centering\small
\caption{Timing and clocks: settling, in-window slope, the CDS offset, and clock feedthrough per volt of swing.}\label{tab:tim}
\begin{adjustbox}{max width=\linewidth}
\begin{tabular}{llrrrrrr}\toprule
setup & channel & reset settling (\textmu s) & window slope (e/\textmu s) & CDS offset (e) & RG (ke/V) & H (ke/V) & V (ke/V)\\\midrule
PA03, 130 K & 101 ch0 & \cellcolor{good}1.09 & \cellcolor{good}4.7 & \cellcolor{good}$-$19 & 6.3 & 25.7 & 49.5\\
 & 101 ch1 & \cellcolor{good}0.95 & \cellcolor{good}5.0 & \cellcolor{good}$-$20 & 5.9 & 26.0 & 52.5\\
 & 101 ch2 & \cellcolor{good}0.90 & \cellcolor{good}4.7 & \cellcolor{good}$-$21 & 6.9 & 26.6 & 53.6\\
 & 101 ch3 & \cellcolor{good}1.14 & \cellcolor{good}5.7 & \cellcolor{good}$-$24 & 5.9 & 26.5 & 50.4\\
 & 102 ch3 & \cellcolor{fair}1.85 & \cellcolor{fair}8.3 & \cellcolor{good}$-$33 & 4.9 & 10.9 & 8.1\\
 & 102 ch0 \textit{(no charge)} & \cellcolor{fair}1.55 & -- & -- & -- & -- & --\\
 & 102 ch1 \textit{(no charge)} & -- & -- & -- & -- & -- & --\\
 & 102 ch2 \textit{(no charge)} & -- & -- & -- & -- & -- & --\\
\midrule
LSM, 180 K & ACM\,106 ch0 & \cellcolor{good}0.52 & \cellcolor{good}$-$0.6 & \cellcolor{good}$-$5 & 6.5 & 12.1 & 15.3\\
 & ACM\,106 ch1 & \cellcolor{good}0.28 & \cellcolor{good}$-$2.0 & \cellcolor{fair}$-$154 & 9.5 & 23.7 & 26.9\\
 & ACM\,106 ch2 & \cellcolor{good}0.31 & \cellcolor{good}3.2 & \cellcolor{fair}$-$178 & 8.9 & 19.8 & 25.1\\
 & ACM\,106 ch3 & \cellcolor{good}0.43 & \cellcolor{good}0.1 & \cellcolor{fair}$-$285 & 6.5 & 10.9 & 15.0\\
 & ACM\,109 ch0 & \cellcolor{good}0.47 & \cellcolor{fair}$-$10.1 & \cellcolor{good}$-$60 & 7.3 & 16.8 & 19.6\\
 & ACM\,109 ch1 & \cellcolor{good}0.32 & \cellcolor{poor}$-$19.5 & \cellcolor{fair}$-$348 & 11.1 & 30.6 & 35.9\\
 & ACM\,109 ch2 & \cellcolor{good}0.34 & \cellcolor{poor}$-$20.7 & \cellcolor{fair}$-$408 & 12.4 & 35.8 & 40.6\\
 & ACM\,109 ch3 & \cellcolor{good}0.37 & \cellcolor{fair}$-$7.5 & \cellcolor{poor}$-$597 & 9.0 & 22.0 & 26.0\\
 & ACM\,110 ch0 & \cellcolor{good}1.04 & \cellcolor{good}$-$0.4 & \cellcolor{fair}$-$206 & 7.7 & 9.7 & 13.2\\
 & ACM\,110 ch1 & \cellcolor{good}0.66 & \cellcolor{good}$-$0.5 & \cellcolor{poor}$-$610 & 6.0 & 8.3 & 12.5\\
 & ACM\,110 ch2 & \cellcolor{good}0.82 & \cellcolor{good}$-$0.1 & \cellcolor{poor}$-$983 & 9.2 & 12.0 & 18.2\\
 & ACM\,110 ch3 & \cellcolor{good}0.94 & \cellcolor{good}$-$0.2 & \cellcolor{poor}$-$1548 & 5.8 & 7.0 & 10.8\\
\bottomrule\end{tabular}\end{adjustbox}
\par\smallskip\footnotesize Grades --- settling: \colorbox{good}{$x\le1.2$} \colorbox{fair}{$x\le2$}; slope: \colorbox{good}{$|x|\le6$} \colorbox{fair}{$|x|\le15$}; offset: \colorbox{good}{$|x|\le100$} \colorbox{fair}{$|x|\le500$}; otherwise \colorbox{poor}{poor}. RG, H and V: relative only, not graded (they include board pickup and the whole transfer waveform).
\end{table}

\begin{itemize}[nosep]
\item The LSM outputs settle faster after reset (0.28--0.52\,\textmu s on boards 106/109) than PA03 ($\sim$1\,\textmu s).
\item \textbf{Board 109's signal windows are not settled} ($-7$ to $-21$\,e/\textmu s), and the LSM CDS offsets are up to
  80$\times$ PA03's ($-150$ to $-1550$\,e): the wait after the summing-well pulse (3.5\,\textmu s) is too short there.
  PA03 drifts $+5$\,e/\textmu s on all channels: worth a longer wait too.
\end{itemize}

% ===================================================================================================================
\clearpage
\section{The setup}
\what{The repeating waveform is removed from the trace (the median skip cycle); what is left is noise.
\textbf{Slow noise 0.2--1\,kHz}: its spectral density at low frequency: pickup and wander from cabling, grounding,
supplies. The board-102 channels without charge (1.5--1.6\,\textmu V/$\sqrt{\text{Hz}}$) are the reference: same electronics
and transistors, no charge. \textbf{Skip corr.}: the correlation
of consecutive skips (slow noise as the skipper sees it; LSM values corrected for their 7-skip pixels).
\textbf{Bursts}: share of 1\,ms blocks with 1.5$\times$ the median noise. \textbf{Static}: the same band in a record with
all clocks frozen (LSM only). \textbf{Blank-image charge}: electrons per pixel in blank images (dark current + spurious
charge; at 180\,K mostly dark current, so not graded on the same footing). \textbf{Setup at 500 skips}: what the
setup actually costs an image. The per-skip $\sigma_\text{CDS}$ cannot see slow pickup (the CDS suppresses
$f<1$\,kHz as $f^2$), but a 500-skip pixel lasts 16\,ms and does not average it away. The static-record spectrum at
100\,Hz--1\,kHz is integrated through the exact CDS + 500-skip filter and divided by the gain: electrons on the image's
scale, and its share of the 500-skip variance.}
\begin{table}[H]\centering\small
\caption{The setup (and the charge in blank images).}\label{tab:setup}
\begin{adjustbox}{max width=\linewidth}
\begin{tabular}{llrrrrrr}\toprule
setup & channel & slow 0.2--1\,kHz (\textmu V/$\sqrt{\text{Hz}}$) & 1--10\,kHz (\textmu V/$\sqrt{\text{Hz}}$) & skip corr. & bursts & static 0.2--1\,kHz & blank-image charge (e/pix)\\\midrule
PA03, 130 K & 101 ch0 & \cellcolor{poor}35.3 & 3.80 & \cellcolor{poor}0.71 & \cellcolor{poor}0.11 & -- & \cellcolor{good}0.20\\
 & 101 ch1 & \cellcolor{poor}36.0 & 3.85 & \cellcolor{poor}0.69 & \cellcolor{poor}0.11 & -- & \cellcolor{good}0.28\\
 & 101 ch2 & \cellcolor{poor}35.0 & 3.74 & \cellcolor{poor}0.68 & \cellcolor{poor}0.08 & -- & \cellcolor{fair}0.49\\
 & 101 ch3 & \cellcolor{poor}34.8 & 3.81 & \cellcolor{poor}0.64 & \cellcolor{poor}0.11 & -- & \cellcolor{fair}0.43\\
 & 102 ch3 & \cellcolor{good}1.5 & 1.11 & \cellcolor{good}0.19 & \cellcolor{good}0.00 & -- & \cellcolor{poor}1.35\\
 & 102 ch0 \textit{(dead)} & \cellcolor{good}1.6 & 1.10 & \cellcolor{fair}0.26 & \cellcolor{good}0.00 & -- & --\\
\midrule
LSM, 180 K & ACM\,106 ch0 & \cellcolor{fair}6.1 & 1.66 & \cellcolor{good}0.03 & \cellcolor{good}0.00 & \cellcolor{poor}3.65 & \cellcolor{poor}1.10\\
 & ACM\,106 ch1 & \cellcolor{fair}12.9 & 2.28 & \cellcolor{good}0.01 & \cellcolor{good}0.00 & \cellcolor{fair}1.26 & \cellcolor{poor}1.15\\
 & ACM\,106 ch2 & \cellcolor{fair}9.5 & 1.79 & \cellcolor{good}$-$0.00 & \cellcolor{good}0.00 & \cellcolor{fair}2.90 & \cellcolor{poor}1.06\\
 & ACM\,106 ch3 & \cellcolor{fair}5.2 & 1.59 & \cellcolor{good}$-$0.01 & \cellcolor{good}0.00 & \cellcolor{fair}1.46 & \cellcolor{poor}1.20\\
 & ACM\,109 ch0 & \cellcolor{fair}7.9 & 1.89 & \cellcolor{good}0.05 & \cellcolor{good}0.00 & \cellcolor{fair}2.07 & \cellcolor{poor}1.02\\
 & ACM\,109 ch1 & \cellcolor{poor}29.2 & 3.89 & \cellcolor{good}0.10 & \cellcolor{poor}0.21 & \cellcolor{fair}2.33 & \cellcolor{poor}0.85\\
 & ACM\,109 ch2 & \cellcolor{poor}15.1 & 2.73 & \cellcolor{good}0.16 & \cellcolor{good}0.00 & \cellcolor{fair}1.13 & \cellcolor{poor}1.02\\
 & ACM\,109 ch3 & \cellcolor{fair}7.3 & 1.75 & \cellcolor{good}0.08 & \cellcolor{good}0.00 & \cellcolor{fair}1.22 & \cellcolor{poor}0.90\\
 & ACM\,110 ch0 & \cellcolor{good}3.6 & 1.53 & \cellcolor{good}0.03 & \cellcolor{fair}0.02 & \cellcolor{good}0.58 & \cellcolor{poor}1.11\\
 & ACM\,110 ch1 & \cellcolor{good}4.2 & 1.45 & \cellcolor{good}0.01 & \cellcolor{good}0.00 & \cellcolor{good}0.53 & \cellcolor{poor}1.12\\
 & ACM\,110 ch2 & \cellcolor{good}2.6 & 1.57 & \cellcolor{good}0.02 & \cellcolor{good}0.00 & \cellcolor{good}0.27 & \cellcolor{poor}0.85\\
 & ACM\,110 ch3 & \cellcolor{poor}19.8 & 2.48 & \cellcolor{fair}0.37 & \cellcolor{poor}0.33 & \cellcolor{good}0.41 & \cellcolor{poor}1.12\\
\bottomrule\end{tabular}\end{adjustbox}
\par\smallskip\footnotesize Grades --- slow noise: \colorbox{good}{$x\le5$} \colorbox{fair}{$x\le15$}; corr.: \colorbox{good}{$x\le0.2$} \colorbox{fair}{$x\le0.5$}; bursts: \colorbox{good}{$x\le0.02$} \colorbox{fair}{$x\le0.08$}; static: \colorbox{good}{$x\le1$} \colorbox{fair}{$x\le3$}; charge: \colorbox{good}{$x\le0.3$} \colorbox{fair}{$x\le0.6$}; otherwise \colorbox{poor}{poor}.
\end{table}

\begin{table}[H]\centering\small
\caption{Setup history of PA03 at equal trace length (0.2\,s): the cabling change seen from the trace alone.}\label{tab:hist}
\begin{tabular}{lrrrrr}\toprule
trace set & CDS noise & slow noise 0.2--1\,kHz & skip-to-skip corr. & bursts & 1--7.5\,MHz\\
 & (nV/$\sqrt{\text{Hz}}$) & (\textmu V/$\sqrt{\text{Hz}}$) & & & (nV/$\sqrt{\text{Hz}}$)\\\midrule
101 ch1, October (before cabling; 0.2\,s traces) & \cellcolor{fair}20.2 & \cellcolor{poor}48.8 & \cellcolor{poor}0.71 & \cellcolor{fair}0.08 & 29.2\\
101 ch1, Nov 2--3 & \cellcolor{fair}19.5 & \cellcolor{poor}42.8 & \cellcolor{poor}0.75 & \cellcolor{fair}0.12 & 29.5\\
101 ch1, Nov 10--12 & \cellcolor{fair}19.6 & \cellcolor{poor}34.7 & \cellcolor{poor}0.68 & \cellcolor{fair}0.11 & 29.4\\
102 ch0, no charge (same-board reference) & \cellcolor{good}17.0 & \cellcolor{good}1.6 & \cellcolor{good}0.26 & \cellcolor{good}0.00 & 10.6\\
\bottomrule\end{tabular}\end{table}

\begin{itemize}[nosep]
\item \textbf{The setup in electrons (LSM, static records):} 0.01--0.02\,e at 500 skips on ACM110 (1--2\,\% of the
  variance), 0.04--0.12\,e on ACM106/109, 0.14\,e on ACM106 ch0 (67\,\%). Channels that look identical per skip (3.3--4.4\,e)
  differ by up to two-thirds of their 500-skip noise. The same static spectra predict the measured per-skip noise of the
  traces to 0.97 (median), so the integration is sound; the images are 1.4$\times$ the predicted total (180\,K dark charge).
\item \textbf{Why only from static records:} the clocked traces carry slow variation that is not image noise. PA03's
  trace CDS series has $\sim$3.5\,e at 128 skips, $\sim$12$\times$ what its images show (0.144\,e at 500 skips); on the
  LSM the clocked-trace 0.2--1\,kHz level is 2--50$\times$ the static one. Clocked-trace numbers (the column
  ``slow 0.2--1\,kHz'' and Table~\ref{tab:hist}) are valid for comparisons made the same way (October vs November, board
  101 vs 102), not as absolute setup noise. PA03 needs static records for the electron number.
\item \textbf{Two concrete board-101 setup items, from one trace:} its video chain recovers with an 8.2\,ms constant
  after every horizontal transfer ($-$1\,800\,ADU; board 102 barely moves), and it rings at 3.1\,MHz after fast edges
  (board 102: 0.55--0.7\,MHz). The recovery drives the within-pixel CDS drift, and the CDS leaks 0.2--0.5\,\% of the
  baseline, which explains much of board 101's trace slow noise. Gain and supplies are stable to 0.3\,\% per skip.
\item \textbf{PA03 board 101 is the noisy setup}: 35\,\textmu V/$\sqrt{\text{Hz}}$ on all four channels, 10\,\% bursty blocks,
  skip correlation 0.64--0.71. The cabling change lowered it by 29\,\% at equal trace length (Table~\ref{tab:hist}); it is still 3--13$\times$
  the LSM. The images tolerate it: board 101's 500-skip images reach the white-noise expectation (Table~\ref{tab:noise}).
\item \textbf{Board 102 does not have it} (ch3 1.5, like ch0): the slow pickup enters on board 101's side
  (its CCD connection, cable or ACM), not through the whole cryostat.
\item LSM: ACM110 ch0--ch2 and ACM106 ch3 are close to that reference. ACM110 ch3 has a 30.6\,kHz comb with
  harmonics (external interference, only while clocking); ACM109 ch1 has bursts.
\end{itemize}
\fig[0.85]{raw_psd_oct_nov.png}{PA03 101 ch1: the trace spectrum before and after the cabling change, with a board-102 channel without charge.}
\fig[0.6]{static_psd.png}{LSM: the static noise records (all clocks frozen) of all twelve channels.}

% ===================================================================================================================
\clearpage
\section{Channel by channel}\label{sec:verdicts}
\begin{itemize}[leftmargin=*]
\item \textbf{PA03 101 ch0, ch1:} the best amplifiers of the study (smallest sense nodes, 0.144--0.147\,e images, low 1/f,
  low trap upper limits); their images are at the white-noise expectation despite board 101's pickup.
\item \textbf{PA03 101 ch2 (excluded since v1):} the transistor is fine on average, but it has the most modulated noise level,
  2.5$\times$ the charge of ch0 in blank images (output-stage glow, spurious charge at the amplifier end of its register),
  and 0.164\,e images. Its defect shows in charge and in episodes, not in the per-skip noise.
\item \textbf{PA03 101 ch3:} the highest per-skip noise of board 101 (4.4\,e) and the highest 1/f corner (39\,kHz); excess
  slow noise in the images (Allan deviation).
\item \textbf{PA03 102 ch3:} live amplifier, quiet setup, but no charge reaches it (mapping to be checked); slow settling.
\item \textbf{ACM106:} clean setup and transistors; images 0.21--0.33\,e (ch1 worst) at 180\,K; CDS offsets up to $-285$\,e.
\item \textbf{ACM109:} windows not settled on all four channels; \textbf{ch1 has intermittent $\sim$1\,ms noise bursts} (one trace;
  its static spectrum is normal); ch3's trap fit fails.
\item \textbf{ACM110:} the quietest setup (ch0--ch2); \textbf{ch1 has more 1/f noise} (174\,kHz corner); \textbf{ch3 picks up
  a 30.6\,kHz interference comb}; the largest CDS offsets ($-600$ to $-1550$\,e).
\end{itemize}
\fig[0.75]{anatomy_ACM109_ch1.png}{ACM109 ch1: the channel with noise bursts. Its whole-trace dissection.}
\fig[0.75]{anatomy_ACM110_ch3.png}{ACM110 ch3: the interference comb (bottom left, lines at 31, 61, 92, 123\,kHz).}

% ===================================================================================================================
\clearpage
\section{What needs a voltage scan (PA03 only so far)}\label{sec:scan}
The LSM data has one voltage setting and no simultaneous strips; these PA03 quantities need a scan or strips:
\begin{table}[H]\centering\small
\begin{tabular}{lll}\toprule
quantity & PA03 board 101 & needs\\\midrule
per-line mutual capacitance to the sense node & RG 0.15--0.37, SW 0.08--0.13, H 0.011--0.016, V $<0.003$\,fF & clock-voltage scan\\
follower gain $A$, $C_\text{sn}$ from the layout inversion & $A\approx0.8$--0.9, $C_\text{sn}\approx80$--90\,fF & scan + layout\\
buried-channel depth & $\approx0.6$\,\textmu m & reset charge injection vs RG\\
channel-potential knees & VR $-2.5$ to $-3.8$\,V & VR/VDD scan\\
oxide traps from the window scan (own noise, shared part removed) & $N_t=1.9$--$2.8\times10^{16}$ & window scan + strips\\
crosstalk, video-line capacitance, shared return & 2--4\,$\Omega$ ch0/ch2 -- ch1/ch3 & simultaneous strips\\
dark-charge response to VSUB & ch0/ch2 weak (register) & VSUB scan\\\bottomrule
\end{tabular}\end{table}

\section*{Caveats}
\begin{itemize}[nosep]
\item Board-102 ch0--ch2 were first treated as ``dead'' (no transistor, electronics only). They have output transistors
  and only receive no charge; their 1/f fits look like transistors (corner 100--130\,kHz, $N_t\approx2\times10^{16}$).
  They are used here as a same-board reference without charge, not as the bare-electronics floor.
\item LSM nV/$\sqrt{\text{Hz}}$ values assume PA03's design chain gain (no independent chain calibration on ACM106/109/110); their electron
  values use their own image gains and are unaffected.
\item LSM image extension EXT$n$ is taken as channel $n-1$ (as on PA03); to be confirmed.
\item LSM: one trace per channel; a single unusual trace can move a row. PA03: 8 traces of valid packets (93 for the p90 columns).
\item Temperatures differ (130 vs 180\,K): thermal noise, 1/f power and dark current all grow with temperature; the
  setup numbers (room-temperature electronics and cabling) do not.
\end{itemize}

\section*{Reproduce}
{\small\begin{verbatim}
.venv/bin/python amp_phys/trace_anatomy.py  --traces <npz|csv> --label <name> --out <dir> [--drop 3]
.venv/bin/python amp_phys/trap_from_trace.py --traces <npz|csv> --label <name> --out <dir> [--t 180]
.venv/bin/python amp_phys/skip_averaging.py --traces <npz|csv> --label <name> --out <dir>
.venv/bin/python amp_phys/static_psd.py --files <psd_*.csv> --out <dir>
.venv/bin/python amp_phys/channel_scorecard.py --out amp_phys/results/per_channel
.venv/bin/python amp_phys/physics_sheet.py --out amp_phys/results/per_channel
.venv/bin/python docs/overleaf_channel_scorecard_2026-11-16/make_tables.py
\end{verbatim}}
\end{document}
